\chapter{The instrument layer: intake, calibration, tare}\label{ch:instrument}

\section{From intensities to $\beta$}

A bisulfite array does not report $\beta$. It reports two fluorescence intensities per probe, methylated $M$ and unmethylated
$U$, in two colour channels, stored in IDAT files. $\beta=M/(M+U+\alpha)$ with a small offset $\alpha$ is computed only after
background correction, dye-bias correction and probe-type handling~\cite{Triche2013,Fortin2017}. This is the methylome's
map-making step, and like map-making it is where most of the instrument's errors enter. The chain does it itself, from the IDATs,
one array at a time (Stage~1, \texttt{chain/\allowbreak{}stage\_1\_idat\_calibration.py}): noob background and dye-bias correction with
probe-type normalisation, using only the array's own control and out-of-band probes. No other array and no reference matrix enters
this step. Stage~1 also records what the array reports about itself: per-probe detection, control-probe intensities and the spread of the
SNP probes.

\section{Intake: the array's own numbers}

Before any cell is read, Stage~0 asks whether the array is a measurement at all (\texttt{chain/\allowbreak{}stage\_0\_intake.py}). It checks
the manifest (array type, declared sex and age, a hashed identifier), the integrity hash of the IDAT pair, the control probes, per-probe
detection, the sample call rate and the sex check. Every quantity is the array's own; none refers to a person. The lines are read from
\texttt{Runtime Matrices/\allowbreak{}Intake/\allowbreak{}intake\_thresholds\_v1.json}:

\begin{center}\small
\begin{tabular}{@{}ll@{}}\toprule
call rate & decision \\\midrule
$\ge0.98$ & proceed \\
$0.93$--$0.98$ & proceed with penalty, flagged \\
$<0.93$ & quarantine: no reading is written \\\bottomrule
\end{tabular}
\end{center}
The call rate is the fraction of probes that are detected and carry at least three beads; a probe is detected when its Gaussian detection $p$,
$1-\Phi\big((I-\mu_{\rm bg})/\sigma_{\rm bg}\big)$ against the array's own negative-control background, is below 0.01
(\texttt{stage\_0\_intake.py}). Stage~1 records a second detection statistic, poobah $p\le0.05$, and a call rate on it; these are recorded, not
gated. The 0.93 line was set from a 48-array measurement across four laboratories: medians 0.985 (minimum 0.979), 0.975 (minimum 0.891) and 0.953
(minimum 0.932) in three, and 0.878 (maximum 0.928) in the fourth (\texttt{doors/\allowbreak{}PROC\_INTAKE\_01\_OUTCOME.md}). \measured{} The line
sits in the gap between the fourth laboratory's best array and the third's worst; one array of the second laboratory
(0.891) falls below it. That measurement used poobah $p\le0.05$, not the statistic the gate applies; the line has not yet been
re-measured on the gate's own statistic. \openprob

Three further rules decide the verdict. A sex mismatch between the array and the manifest quarantines the specimen and is never silently
corrected. A deferred detection or call-rate check is a hard quarantine: a check that cannot run never advances a specimen
(PROC-INTAKE-01, bar B5). Bead counts are extracted at the Stage~0 hand-off; a low-bead warning gives \emph{proceed with penalty}, and probes with fewer than three
beads do not count toward the call rate. The bisulfite-conversion threshold is recorded and not gated: as written it refused all 731 healthy
arrays it was tried on, and it is printed as provisional. \openprob{} A quarantined
specimen produces no reading; the run stops with exit status 2.

\paragraph{Probes at background.} A probe at background reads both channels near noise, so its $\beta$ sits near 0.33--0.42 whatever the
true methylation (PROC-INTAKE-01, diagnostic after bar B3). \measured{} Met-A identity sites sit at 0.75--0.95 or 0.05--0.25
(Chapter~\ref{ch:identity}), where such a value raises $H(\beta)$ toward one bit. Stage~1 therefore removes every probe whose detection $p$
exceeds 0.05 before any stage reads a $\beta$.

\section{One platform per reference}

The chain reads EPIC v1 arrays only. It refuses a specimen when Stage~0 reports another array type, when probe names carry the EPIC v2
design suffix, or when the $\beta$ vector holds 700,000 probes or fewer (\texttt{conductor\_v3.py}, \texttt{platform\_refusal}). The reason is
measured: a reference built on one platform does not transfer to another (Chapter~\ref{ch:meta}). A 450K neutrophil reference is not
yet frozen. \openprob

\section{The noise index}

Array noise raises $H$ at every site, and the identity sites are where the gauge reads it. The chain therefore measures the noise of each
array on sites that biology holds still. The noise index $N$ is the mean $H(\beta)$ over 48,528 EPIC sites at which every purified blood
group of the reference has mean $\beta\le0.03$ (40,882 sites) or $\ge0.97$ (7,646 sites) with every group's standard deviation $\le0.02$; no
neutrophil identity site is among them (\texttt{noise\_sites\_EPIC\_v1.json}; \texttt{doors/\allowbreak{}DEV\_NOISE\_01\_OUTCOME.md}). $N$ is
computed when at least 90~\% of those sites are measured. \calibrated{} (the site list)

\begin{center}\small
\begin{tabular}{@{}lrll@{}}\toprule
arrays & $n$ & $N$, median (range) & Met-A, median \\\midrule
reference neutrophils & 12 & 0.1285 (0.1223--0.1489) & 1.001 \\
second laboratory, donor 1 & 24 & 0.1504 (0.1067--0.1655) & 1.069 \\
second laboratory, donor 2 & 24 & 0.1656 (0.1365--0.2434) & 1.192 \\\bottomrule
\end{tabular}
\end{center}
\measured{} (\texttt{DEV\_NOISE\_01\_OUTCOME.md}; the 12 reference rows are the six physical arrays deposited twice.) Within the second laboratory
Met-A follows $N$ ($\rho=0.79$ and 0.83). $N$ does not explain all of it: some second-laboratory arrays with $N$ at the reference level still read
about 1.07. That residual is the reason the chain also tares every reading against references run beside it.

\section{The tare: references run beside the specimen}

A physical instrument is zeroed against a known input measured the same way. Chain v3 does this with healthy reference specimens of the same
type, run on the same slide, else in the same batch, through the same chain (Stage~T, \texttt{conductor\_v3.py}, \texttt{stage\_t\_tare}). The
references are part of the run, not a stored population, and the specimen is never among its own references.
\begin{itemize}
\item \textbf{Median tare} (at least 3 references): $A_{\rm rel}=A/\mathrm{median}(A_{\rm ref})$. There is no fitted correction.
\item \textbf{Noise gate:} the noise index $N$ (mean $H(\beta)$ over the 48,528 sites every purified blood group holds fixed) is measured on
      every array and gates the gauge state: outside the range of the purified reference arrays the state is withheld.
\item \textbf{Detection limit:} $2\times$ the references' spread divided by the shift that a 1~\% loss of the neutrophil pattern would cause at this
      specimen's own fraction. Every tared reading states the smallest loss it could have shown.
\end{itemize}
With fewer than three references the reading is untared. Isolated neutrophils are then drawn on the gauge against their own reference,
labelled \emph{untared}, with the own-reference state kept in the state text; whole blood gets no gauge position until it is tared
(\texttt{conductor\_v3.py}, \texttt{report\_v3.py}). The operating procedure is in the operations manual
(\texttt{Biological\_Physics/\allowbreak{}MethylPhys/\allowbreak{}manual/\allowbreak{}OM\_v3/\allowbreak{}OM\_v3.tex}).

The acceptance run shows what the tare does and what it has not yet been tested on (\texttt{chain\_tests/\allowbreak{}CHAIN\_V3\_ACCEPTANCE\_RUN3.md},
\texttt{chain\_acceptance.csv}). Untared whole-blood Met-A read 0.943--0.968 on six DNA mixtures from the reference laboratory and 1.073--1.115 on
five whole bloods from another laboratory. The acceptance script re-reads whole blood only, with a median tare whose references are the other
specimens of the same group: for the second laboratory these were the other remission bloods, not healthy references. So tared, they read
0.987--1.021 and 0.986--1.032, all in Normal. \measured{} That is a test of the arithmetic, not of the tare against healthy references; isolated
specimens were not tared in that run. Isolated neutrophils need the tare as well: a second laboratory's purified neutrophils read 0.86--1.26
against the frozen reference with no tare, and the spread followed the arrays' noise, not purity
(\texttt{PROC\_NEUT\_TEST\_01\_T2\_OUTCOME.md}). \measured

In one laboratory's 76 healthy-control whole-blood arrays, untared $A$ tracks the noise index and the neutrophil fraction: a linear fit
$A\approx0.60+0.205\,f_{\rm neu}+2.51\,N$ has $R^2=0.85$ (\texttt{DEV\_NOISE\_02\_OUTCOME.md}). \fitted{} The fit was made after looking at
those data and the chain does not apply it.

\paragraph{A tare from the SNP probes.} Every array carries probes on common polymorphisms that read $\beta=0$, $\tfrac12$ or $1$ by genotype. A
linear tare fitted to them was tested on 768 arrays from four laboratories and not commissioned: it moved the reading by a near-constant amount on
94~\% of arrays, its scale term varied more within a chip than between chips, and neither parameter predicted the reading ($r=-0.14$ and $-0.07$)
(\texttt{PROC\_TARE\_01\_OUTCOME.md}). \measured{} Compression at $\beta=0$ and 1 says little about the response where identity sites sit.

\section{What remains open}

The planned tare is a slide tare: fully methylated and fully unmethylated control DNA run on every slide, so that each slide is zeroed
against known inputs measured the same way. It is \openprob{} until a laboratory runs it. A second route, a nonlinear response model using the
control probes as well as the SNP clusters, is not started. The intake call-rate line needs re-measuring on the gate's own detection statistic,
the bisulfite-conversion line needs calibration, and a 450K reference needs freezing before 450K arrays can be read. \openprob
